% Generated by scripts/make_table.py from tables/evidence_matrix.csv. Do not edit by hand.
\begin{landscape}
\begingroup
\footnotesize
\singlespacing
\setlength{\tabcolsep}{4pt}
\begin{longtable}{>{\raggedright\arraybackslash}p{1.0in}>{\raggedright\arraybackslash}p{1.1in}>{\raggedright\arraybackslash}p{0.9in}>{\raggedright\arraybackslash}p{1.7in}>{\raggedright\arraybackslash}p{1.3in}>{\raggedright\arraybackslash}p{1.9in}>{\raggedright\arraybackslash}p{0.3in}}
\caption{Comparative evidence matrix: experimental studies of rodent receivers cited in this review, in chronological order. Setting: lab, field, or semi (wild-caught animals tested in captivity or semi-natural conditions). nr, not reported in the sources consulted for this draft. Q, review question(s) addressed. The full matrix, including limitations and verification status, is in tables/evidence\_matrix.csv.}\label{tab:evidence} \\
\toprule
\textbf{Study} & \textbf{Species (setting)} & \textbf{Receiver experience} & \textbf{Stimuli; controls} & \textbf{Endpoints} & \textbf{Key result} & \textbf{Q} \\
\midrule
\endfirsthead
\toprule
\textbf{Study} & \textbf{Species (setting)} & \textbf{Receiver experience} & \textbf{Stimuli; controls} & \textbf{Endpoints} & \textbf{Key result} & \textbf{Q} \\
\midrule
\endhead
\bottomrule
\endlastfoot
\citet{Davis1970} & Rat (Sprague-Dawley, male) (lab) & Naive & Startle tones; Exp I: 2-s vs 16-s intervals (n=72); Exp II: regular vs variable at 8-s mean (n=144); Schedule comparison & Startle amplitude; retest at 1 min and 24 h & 2-s spacing gave more decrement in training but 16-s spacing gave lower retest startle; regular timing gave more habituation than variable at matched mean & 6, 7 \\ \addlinespace[3pt]
\citet{EvansHammond1983a} & Rat (albino, male) (lab) & Laboratory-reared & Recorded distress squeal at 80 or 100 dB; Exp 1: 10 trials 3 min apart (n=8 per group); Exp 2: retention after 0, 1, or 7 days (n=120); Exp 3: daily trials (n=32); Synthetic mimic matched in burst timing & Orienting (lick suppression) & Natural squeal exceeded mimic at 100 but not 80 dB; all habituated rapidly; the 100-dB squeal recovered more after 1- and 7-day rests and habituated more slowly with 24-h spacing & 1, 6, 7 \\ \addlinespace[3pt]
\citet{EvansHammond1983b} & Rat (lab) & Laboratory-reared & Natural distress squeal and mimic across contexts; Synthetic mimic & Orienting; generalization of habituation & Generalization of habituation across contexts depended on stimulus significance & 6 \\ \addlinespace[3pt]
\citet{HendrieNeill1991} & Mouse (DBA/2, male) (lab) & Laboratory-bred; predator-naive & 2-min calls of 10 species (predators: barn, tawny, little owl, kestrel, fox), 60-70 dB, after 3 min of human voice; n=10 per group; Human voice; non-predator calls (mallard, gull, oystercatcher, great tit, red deer) & Cage-, self-, and call-oriented behavior; defense; food intake & Barn and tawny owl calls increased self-oriented, call-oriented, and defensive behavior; kestrel and gull calls also increased defense; tawny owl call reduced feeding by about half (t(18)=2.67) & 2, 4 \\ \addlinespace[3pt]
\citet{Mateo1996} & Belding's ground squirrel (semi) & Juveniles; early auditory experience & Alarm calls (whistles, trills); nr & Heart rate; behavior & Heart-rate discrimination of alarm calls preceded behavioral discrimination & 2 \\ \addlinespace[3pt]
\citet{BlumsteinArmitage1997} & Yellow-bellied marmot (field) & Free-living & Alarm calls varying in type, rate, and amplitude; nr & Vigilance and behavior & Calls tracked risk rather than predator type; receivers used call rate as risk information & 1, 4, 7 \\ \addlinespace[3pt]
\citet{Hare1998} & Richardson's ground squirrel (juvenile females) (field) & Free-living juveniles & Five calls about 3 min apart, each a repeated call of five chirps from a juvenile female, 84-91 dB; four of one exemplar, then a new exemplar from the same caller or a call from a different caller; Same caller, new exemplar & Vigilance duration; posture & Vigilance fell by the second call; a different caller restored it, a new exemplar from the same caller did not; neighbours' calls evoked longer vigilance than strangers' (trend, P=0.06) & 5, 6 \\ \addlinespace[3pt]
\citet{HendrieEtAl1998} & Voles, dormice, jerboas, spiny mice (semi) & Wild-caught & Owl calls; nr & Behavior & Most species changed behavior; spiny mice responded less & 2 \\ \addlinespace[3pt]
\citet{Shriner1998} & Yellow-bellied marmot; golden-mantled ground squirrel (field) & Free-living & Conspecific and heterospecific alarm calls; Sympatric bird song & Vigilance & Both species distinguished alarms from song and responded similarly to conspecific and heterospecific alarms & 2 \\ \addlinespace[3pt]
\citet{EilamEtAl1999} & Gunther's vole; spiny mice (semi) & Wild-caught & Owl calls; nr & Freezing or flight; corticosterone or cortisol & Voles froze or fled and raised corticosterone; spiny mice raised cortisol without behavioral change & 2, 3, 4 \\ \addlinespace[3pt]
\citet{BaackSwitzer2000} & Eastern chipmunk (field) & Free-living adults & Conspecific alarm calls; Control noise; silence & Vigilance; run to cover; re-emergence latency; seed intake & Alarms increased vigilance, shortened the run to cover, delayed re-emergence, and reduced food intake & 4 \\ \addlinespace[3pt]
\citet{HareAtkins2001} & Richardson's ground squirrel (juvenile females) (field) & Free-living juveniles & Ten exemplars each of a reliable caller (paired with a stuffed badger except on playbacks 5 and 10) and an unreliable caller (no predator), 84-91 dB; Within-subject reliable vs unreliable caller & Postural vigilance; vigilance duration & Postural responses to the unreliable caller declined and differed from the reliable caller by playback 5 (P=.014) and 10 (P=.007); vigilance duration showed the same pattern only as a trend (P=.075) & 5 \\ \addlinespace[3pt]
\citet{MongeauEtAl2003} & Mouse (lab) & With or without prior footshock & 17-20 kHz FM sweeps, 85 dB; Home vs unfamiliar cage; prior shock & Flight; freezing; neural activity & Cage familiarity and prior shock shifted responses between flight and freezing & 1, 4 \\ \addlinespace[3pt]
\citet{BlumsteinDaniel2004} & Yellow-bellied marmot (field) & Free-living & Habituation-recovery alarm playback; juvenile vs adult-female callers; familiar vs unfamiliar callers; Caller switch & Foraging suppression & Habituation to one caller recovered with a new caller; juvenile calls more evocative; familiarity null (Z=-0.30, p=.767) & 1, 5, 6 \\ \addlinespace[3pt]
\citet{BlumsteinEtAl2004} & Yellow-bellied marmot (field) & Free-living; reliability manipulated & Alarms from reliable vs unreliable callers; degraded (distant) vs intact calls; one vs two callers; Baseline foraging; playback order & Foraging suppression; vigilance & Unreliable, distant-sounding, and two-caller playbacks were more evocative & 5 \\ \addlinespace[3pt]
\citet{SloanEtAl2005} & Richardson's ground squirrel (field) & Free-living & Chirps, whistles, chucks; nr & Vigilance & Call types carried different information and altered vigilance differently & 1, 7 \\ \addlinespace[3pt]
\citet{KiriazisSlobodchikoff2006} & Gunnison's prairie dog (field) & Free-living (two colonies, 40-50 animals each) & 11-s alarm-call bouts originally elicited by humans, dogs, or coyotes; 45 dB at 50 m; Rubber-toy sound matched in rhythm, frequency range, and amplitude & First and second escape responses & Playbacks reproduced the escape responses to the corresponding live predators: burrow diving (humans), running to burrow and standing at rim (coyotes), alert posture in place (dogs) & 1, 4 \\ \addlinespace[3pt]
\citet{SloanHare2006} & Richardson's ground squirrel (field) & Free-living adults & Juvenile alarm calls with manipulated rate and bout length; Unmanipulated calls & Vigilance & Adults ignored rate and bout-length changes in juvenile calls & 7 \\ \addlinespace[3pt]
\citet{EndresEtAl2007} & Rat (Sprague-Dawley, male; 59 total) (lab) & Naive, then conditioned & Natural 22-kHz calls; 22-kHz sine tones; 22-kHz-range noise; other ultrasonic stimuli; Other frequency bands & Defensive behavior; freezing; conditioning and extinction & No single stimulus altered defense in naive rats (n=9); pooled 22-kHz range increased freezing; learned 22-kHz fear unusually persistent & 2 \\ \addlinespace[3pt]
\citet{HerryEtAl2007} & Mouse (and human) (lab) & Naive & 5-kHz, 20-ms, 85-dB pulses at 5 Hz (SD 140 ms) or 0.7 Hz (mean 1429 ms, SD 1368 ms); Predictable vs jittered sequences & Basolateral amygdala units; c-Fos; avoidance & Jittered timing prevented rapid amygdala habituation, increased c-Fos, and promoted avoidance & 7 \\ \addlinespace[3pt]
\citet{WohrSchwarting2007} & Rat (Wistar, male) (lab) & Juvenile and adult cohorts & Natural 50-kHz calls; sine-wave approximations; 22-kHz calls; Background noise & Approach & Natural and synthetic 50-kHz calls both induced approach, especially in adults & 1 \\ \addlinespace[3pt]
\citet{SadanandaEtAl2008} & Rat (Wistar, male) (lab) & Naive & 22-kHz calls; flat and FM 50-kHz calls; Background noise; silence & Approach; c-Fos & 50-kHz: approach and c-Fos in frontal cortex, accumbens, thalamus; 22-kHz: perirhinal cortex, amygdala, periaqueductal gray & 1, 3 \\ \addlinespace[3pt]
\citet{SchmidtEtAl2008} & Eastern chipmunk (field) & Free-living & Tufted titmouse mobbing calls; hawk calls; nr & Giving-up density & Foraging changed more to titmouse mobbing calls than to hawk calls & 1, 2 \\ \addlinespace[3pt]
\citet{SwanHare2008} & Richardson's ground squirrel (field) & Free-living & Repeated calls; primary vs later syllables; nr & Vigilance & Primary syllables had especially strong alerting effects & 1, 7 \\ \addlinespace[3pt]
\citet{BlumsteinRecapet2009} & Yellow-bellied marmot (field) & Free-living & Alarm calls with added white-noise nonlinearities; Normal calls; silence-added calls & Foraging & Foraging fell more after calls with added nonlinearities & 1, 4, 6 \\ \addlinespace[3pt]
\citet{KindermannEtAl2009} & Mouse (C57BL/6J), rat (Sprague-Dawley), pallid gerbil; all male (lab) & Laboratory-reared, predator-naive & 3-min calls of owls and raptors, 70-80 dB peak, 0.3-9 kHz; Exp 2-3: calls as conditioned stimuli with footshock (10 pairings or 1); Non-predator bird calls (chaffinch, blackbird, mallard, quail) & Freezing; exploration; vigilant rearing; conditioned freezing & No unconditioned defense to predator calls in any species (power to detect 12\% change in freezing); all calls conditionable after 10 pairings with no predator advantage; one-trial learning only for tawny owl, goshawk, and mallard & 2 \\ \addlinespace[3pt]
\citet{KimEtAl2010} & Rat (male pairs) (lab) & Naive vs footshock-experienced receivers & 22-kHz calls from a fear-conditioned partner; Naive receivers; auditory thalamus lesion or inactivation & Freezing & Only experienced receivers froze; auditory thalamus required & 2, 5 \\ \addlinespace[3pt]
\citet{Mateo2010} & Belding's ground squirrel (semi) & nr & Trills (terrestrial predators) and whistles (aerial predators); nr & Cortisol; heart rate & Trills raised cortisol; whistles were linked to slowed heart rate & 2, 4 \\ \addlinespace[3pt]
\citet{AschemeierMaher2011} & Woodchuck; eastern chipmunk (field) & Free-living & Conspecific and heterospecific alarm calls; nr & Vigilance & Woodchucks increased vigilance to both, more to conspecific alarms & 2 \\ \addlinespace[3pt]
\citet{HareWarkentin2012} & Richardson's ground squirrel (field) & Free-living juveniles & Mother's vs other female's calls; colony vs non-colony callers; nr & Vigilance & No differential response to mother's or colony members' calls & 5 \\ \addlinespace[3pt]
\citet{ParsanaEtAl2012a} & Rat (Sprague-Dawley) (lab) & Naive & 22-kHz and 50-kHz calls; Frequency-matched tones & Freezing; 127 amygdala units & No freezing; matched tones reproduced most firing differences between call types & 1, 2, 3 \\ \addlinespace[3pt]
\citet{ParsanaEtAl2012b} & Rat (lab) & Footshock-experienced & 22-kHz calls; nr & Freezing & Proposed that rats learn from their own 22-kHz calls (autoconditioning) & 2 \\ \addlinespace[3pt]
\citet{GetschowEtAl2013} & Eastern gray squirrel (field) & Free-living & Conspecific alarm; robin (sympatric), blackbird (allopatric, similar), honeyeater (allopatric, dissimilar) alarms; Avian control sounds including robin song & Vigilance; latency & Response varied with playback ($\chi^2$=19.68, p=.0014); robin and blackbird alarms like squirrel alarms; honeyeater less & 2 \\ \addlinespace[3pt]
\citet{WilluhnEtAl2014} & Rat (male, food-restricted) (lab) & Individually housed & Natural 50-kHz and 22-kHz calls, $\sim$69 dB, 15 s, 1-2 min intervals; Time- and amplitude-matched noise; background & Approach; accumbens dopamine (voltammetry) & 50-kHz calls increased approach and dopamine (n=12), both declining rapidly; 22-kHz and noise did not & 1, 3, 6 \\ \addlinespace[3pt]
\citet{JayneEtAl2015} & Eastern gray squirrel (field) & Free-living & Predator calls during foraging; nr & Vigilance; foraging; distance to refuge & Increased vigilance, suppressed and delayed foraging; distance to refuge unchanged & 4 \\ \addlinespace[3pt]
\citet{GeisslerEtAl2016} & House mouse (lab) & Experience with sound meaning & Sounds with context-specific meaning; nr & Auditory cortex activation & Cortical activation depended on the learned meaning of sounds & 3 \\ \addlinespace[3pt]
\citet{CalubEtAl2018} & Rat (lab) & Footshock; laryngeal nerve transection vs sham & 22-kHz playback in a novel context; Sham surgery; vocalized vs not & Freezing & Preventing self-vocalization did not reduce acquired call-evoked freezing & 2 \\ \addlinespace[3pt]
\citet{LillyEtAl2019} & Eastern gray squirrel (field) & Free-living & Hawk call followed by bird chatter; Ambient noise & Vigilance & Vigilance relaxed faster with bird chatter & 2, 5 \\ \addlinespace[3pt]
\citet{FreemanEtAl2020} & Richardson's ground squirrel (females, n=29) (semi) & Laboratory chamber & Chirps or whistles from unfamiliar callers; 49 syllables over 2 min; 80 dB at 1 m; No-call control & c-Fos & Chirps reduced c-Fos in medial amygdala and superior colliculus; most regions unchanged & 3 \\ \addlinespace[3pt]
\citet{OlszynskiEtAl2020} & Rat (Wistar) (lab) & Paired (n=30) vs single (n=31) housing & Natural 22-kHz and 50-kHz calls; Computer-generated tones & Vocalization; heart rate; locomotion & Natural calls usually produced larger effects, with similar directions; single housing enlarged several effects & 1, 4 \\ \addlinespace[3pt]
\citet{SchiavoEtAl2020} & Mouse (female) (lab) & Before vs after maternal experience & Pup-call temporal morphs (natural interval $\sim$175 ms); Synthetic interval shifts & Pup retrieval; auditory cortex & Experience broadened and shifted behavioral and cortical tuning & 1, 2 \\ \addlinespace[3pt]
\citet{TasakaEtAl2020} & Mouse (lab) & Maternal experience & Pup calls; nr & Temporal association cortex activity; retrieval & Temporal association cortex contributes to auditory-driven maternal plasticity & 2, 3 \\ \addlinespace[3pt]
\citet{WatermanMai2020} & Cape ground squirrel (field) & Free-living & Crowned lapwing alarms; conspecific alarms; nr & Vigilance & As vigilant to lapwing alarms as to conspecific alarms & 2 \\ \addlinespace[3pt]
\citet{HernandezEtAl2021} & Abrothrix; Oligoryzomys; Rattus rattus (field) & Free-living & Calls of two owl species; No owl call & Feeding; locomotion; vigilance (camera traps) & Predator- and species-specific changes; moonlight and cover also mattered & 2, 4 \\ \addlinespace[3pt]
\citet{VielleEtAl2021} & Rat (lab) & Cagemates and strangers & Self-administered playback of cagemate vs stranger 50-kHz calls; Subthalamic lesion vs sham & Playback choice; cocaine intake & Preferred stranger calls; only stranger calls reduced cocaine intake; lesions abolished both & 1, 5 \\ \addlinespace[3pt]
\citet{BerzEtAl2022} & Rat (lab) & Playback followed by social interaction & 50-kHz calls; nr & Approach; brain activation & Approach habituated despite contingent social interaction & 5, 6 \\ \addlinespace[3pt]
\citet{GoedhoopEtAl2022} & Rat (Long-Evans, male) (lab) & Conditioned & White noise 70-96 dB (typically 90 dB, 6 s) predicted by a 5-s cue; Intensity; probability & Accumbens core dopamine (voltammetry, n=21); avoidance; foraging; operant choice & Dopamine fell to noise and cue; did not scale with intensity (p=.88); unexpected noise slowed recovery & 3, 7 \\ \addlinespace[3pt]
\citet{WarrenEtAl2025} & Prairie vole (female) (lab) & Pair-bonded & Partner vs stranger male calls, including calls directed at another female; nr & Preference behavior & Discriminated partner from stranger calls & 1, 5 \\ \addlinespace[3pt]
\citet{FuruyamaEtAl2026} & Mouse (C57BL/6 and DAT-Cre, both sexes) (lab) & Head-fixed & Noise bursts 45-105 dB, 10 s; ramped and damped noise; Intensity series & Accumbens and tail-striatum dopamine (photometry); striatal neuron calcium; pupil; place avoidance at 85 dB & Accumbens dopamine fell more with louder noise (9 decrease-type mice); 6 others also showed an onset increase on the first burst; tail-striatum dopamine rose & 3 \\ \addlinespace[3pt]
\end{longtable}
\endgroup
\end{landscape}
